\section{Application to Flock}
\label{sec:spin-flock}

We integrate SPIN--Brakedown into Flock~\cite{flock26} to measure its
effect inside a larger prover. Flock checks batches of Boolean constraints
and reduces them to claims about a committed witness. Our integration
replaces its Ligerito PCS~\cite{ligerito25} with the construction in
Section~\ref{sec:spin-pcs}. The final comparison improves prover throughput
by $1.93\times$ and $1.38\times$ at the two measured workloads.

\paragraph{Adapting the opening interface.}
Flock's final claims include interpolation weights on a small group of
witness coordinates. We incorporate these weights into the PCS row
combination, so the opening proves the functional that Flock actually
requests. A compatible packed representation passes the witness between
the two implementations without repacking it. The current adapter uses
two separate openings with $2110$ column samples each and $K=2^{16}$.
Their conditional fixed-matrix error bounds sum to less than $2^{-102}$.
Composition with Flock's adaptive claims requires the corresponding
extraction argument. The benchmark proves the committed batch's
compression constraints; binding public hash inputs and outputs is a
separate application wrapper.

\paragraph{Improving the surrounding prover.}
During integration, we found that field arithmetic, repeated witness
passes, and temporary storage could obscure the PCS improvement. We
optimized Flock's x86 kernels, reused workspaces, and evaluated the
fixed linear constraints directly. The Ligerito configuration receives
the applicable improvements as well. In particular, its existing
column buffers can supply an intermediate zerocheck value, avoiding
duplicate work. These changes preserve the checked protocol messages.
The comparison below uses both PCS backends within this optimized
implementation, rather than attributing shared improvements to SPIN.

\begin{table}[tb]
\centering
\caption{Flock proofs of BLAKE3-compression constraints on one Ryzen
core. Both configurations use the applicable shared optimizations.
The Ligerito configuration is Fast100.}
\label{tab:spin-flock}
\small
\begin{tabular}{@{}rlrrr@{}}
\toprule
Compressions & PCS & Prover (ms) & Verify (ms) & Proof (MiB) \\
\midrule
% Generated by build_application_tables.py; data SHA-256 27480f2ae253a94829a8b2184e483405bdaf91be3971e18ffaf924e23df35ce3
$16,384$ & SPIN--Brakedown & 111.99 & 38.31 & 6.713 \\
$16,384$ & Ligerito & 216.48 & 27.59 & 0.283 \\
$65,536$ & SPIN--Brakedown & 412.84 & 45.90 & 12.834 \\
$65,536$ & Ligerito & 567.93 & 27.73 & 0.396 \\
\bottomrule

\end{tabular}
\end{table}

\paragraph{Measurement method.}
Both configurations run in one executable with the same deterministic
input stream and compiler settings. They use the single-worker hardware
environment of the standalone PCS experiment. Each workload has four processes per backend
in a balanced serial order. Each process excludes five warmups and
measures four trials; entries average the four process medians.
Prover time includes witness generation, commitment, and proving. Setup,
verification, and outer proof serialization are outside the interval.
Serialization inside the SPIN opening remains timed because its bytes
enter the protocol transcript. All $252$ proofs in the final experiment,
including warmups and reference controls, verify.

\paragraph{Prover costs.}
At the larger workload, commitment takes $100.97$\,ms with SPIN--Brakedown
and $161.00$\,ms with Ligerito. Opening takes $61.17$ and $140.35$\,ms,
respectively. These phases account for about $139$\,ms of the $155$\,ms total
difference. Witness generation is faster in the Ligerito configuration,
and the remaining checking and preparation costs differ with the layouts.
Thus the table measures the complete PCS integrations; it does not isolate
the encoder alone.

The current SPIN--Brakedown integration favors prover time at a substantial
communication cost. At the larger workload, its proof is $12.83$\,MiB,
compared with $0.396$\,MiB for Ligerito, which also verifies faster.
The application result complements the standalone PCS measurements:
SPIN's encoding performance supports a fast commitment path, and that
advantage remains useful after integration into Flock.
